\section{Introduction}

Multi-agent systems built on large language model (LLM) backends exhibit a failure mode that has received insufficient systematic treatment in the literature: \textbf{autonomous drift}. In these systems, agents spawn child agents to delegate sub-tasks, and those children may in turn spawn their own descendants. When this spawning chain proceeds without an enforceable mechanism linking each agent to its originating root, the resulting descendant graph diverges from the operator's intent. Individual agents may continue processing long after their parent has concluded, consume resources on goals that are no longer relevant, or propagate context in ways that corrupt downstream reasoning. We call this phenomenon \emph{drift}, and it represents a fundamental tension between the autonomy that makes agentic systems powerful and the traceability that production deployments require.

The drift problem is not merely an operational inconvenience; it is a structural consequence of unbounded spawning. An operator initiates an agent to accomplish a task. That agent spawns a child to handle a specialized sub-problem. The child spawns a grandchild. Three levels deep, the original agent has already returned control to the operator, but the grandchild has no scheduled termination and no native awareness that its ancestral context has dissolved. If the grandchild encounters an obstacle, it cannot query its parent—because the parent's session has closed. It may retry, spawn further descendants, or persist in a loop. The operator, observing only the root session, has no visibility into the descendant tree and no mechanism to terminate it. This is drift.

Existing approaches to bounding agent behavior are well-known and widely deployed, yet each fails to address drift at its root. Time-to-live (\textbf{TTL}) constraints are the most common remedy: an agent is permitted a fixed budget of processing time, after which its session terminates regardless of task state. While TTL prevents unbounded resource consumption, it does so by destroying agents mid-task, creating orphaned sub-tasks with no provision for graceful completion or context inheritance. A spawned child processing a payment transaction, for instance, cannot be safely interrupted by a TTL; partial state may persist at both the child and any downstream services it has already contacted. Depth limits offer a second approach: the spawning chain is permitted at most $k$ generations, after which further spawning is prohibited. Depth limits impose a hard ceiling on the descendant graph, but they do so without any notion of parent lifecycle. A depth-limited grandchild still has no awareness of whether its parent remains active, and the operator still cannot observe or manage the tree that has formed. Neither TTL nor depth limits address the core issue: the absence of a persistent, immutable linkage between each agent and its root ancestor.

We identify a simple but powerful structural invariant that eliminates drift at its source: the \textbf{immutable parent pointer chain}. Under this model, every agent stores a readonly reference to the agent that spawned it, forming a chain from any descendant back to the root. Critically, this chain is immutable—once established, it cannot be modified, severed, or reparented. The root agent holds a null parent pointer by definition. When an agent spawns a child, the child's parent pointer is locked to the spawning agent at the moment of creation. The child cannot later reparent itself to a different ancestor, and no external authority can reassign the chain. This immutability is the key insight: it makes the ancestor graph a directed tree with no possibility of edge mutation after creation. Because the chain cannot be broken, the operator can always traverse parent pointers from any active agent back to its root, regardless of how many sessions have opened or closed in the intervening generations. Drift is not mitigated or contained—it becomes structurally impossible, because the structural conditions that permit drift (reparenting, orphaned ancestry, unobservable descendant trees) cannot arise when parent pointers are immutable and every agent retains a live reference to its chain.

We formalize drift as follows. Let $G = (V, E)$ be a directed graph where each vertex $v \in V$ represents an agent and each directed edge $(p, c) \in E$ represents a spawning relationship: parent $p$ spawned child $c$. Drift is said to occur when there exists a vertex $v \in V$ such that (1) the session of the root ancestor $r$ of $v$ has terminated, (2) $v$ remains active, and (3) there exists no accessible reference from $v$ to $r$ that permits termination signaling or context inheritance. Under this definition, drift is directly enabled by the mutability or absence of parent pointers. When parent pointers are immutable and retained by every agent for its entire lifetime, condition (3) is unsatisfiable: each agent retains a live chain to its root, making termination signaling and context inheritance always accessible. The system therefore satisfies the drift-free invariant: every active agent is reachable from some live root via its parent chain.

This paper presents \textbf{Recursive Spawning}, a model for multi-agent delegation built on the immutable parent pointer chain as a first-class architectural primitive. We describe the AgentOS runtime that enforces immutability at spawn time, the cascade backstop mechanism that propagates termination signals along the parent chain, and the observable descendant graph that gives operators real-time visibility into all active agents and their lineage. We evaluate the approach against TTL and depth-limited baselines in a simulated production workload, demonstrating zero drift events under the immutable chain model versus measurable drift rates under each prior approach. The remainder of this paper is organized as follows. Section~\ref{sec:background} reviews related work in multi-agent coordination and resource bounding. Section~\ref{sec:model}} defines the recursive spawning model and the drift-free invariant formally. Section~\ref{sec:architecture} describes the AgentOS implementation. Section~\ref{sec:evaluation} presents experimental results. Section~\ref{sec:discussion}} discusses limitations and practical considerations for deployment. Section~\ref{sec:conclusion}} concludes.
